\documentclass[12pt,a4paper]{article}

% -------------------------
% PAQUETES
% -------------------------
\usepackage[spanish]{babel}
\usepackage[utf8]{inputenc}
\usepackage[T1]{fontenc}
\usepackage{times}
\usepackage{graphicx}
\usepackage{setspace}
\usepackage{geometry}
\usepackage{titlesec}
\usepackage{pdflscape}
\usepackage[colorlinks=true,linkcolor=black,urlcolor=blue,citecolor=black]{hyperref}
\usepackage{fancybox}
\usepackage{array}
\usepackage{caption}
\usepackage{float}
\usepackage{tcolorbox}
\tcbuselibrary{listingsutf8}
\usepackage{listings}
\tcbuselibrary{verbatim}

\geometry{left=3cm,right=2.5cm,top=3cm,bottom=3cm}
\setstretch{1.5}

\newcommand{\imgstd}[1]{
    \centering
    \includegraphics[width=0.85\textwidth,height=0.45\textheight,keepaspectratio]{#1}
}

\begin{document}

% -------------------------
% PORTADA
% -------------------------
\begin{titlepage}
\centering
\setstretch{1.3}

{\large \textbf{UNIVERSIDAD NACIONAL DEL ALTIPLANO}}\\[0.3cm]
{\large \textbf{FACULTAD DE INGENIERÍA MECÁNICA ELÉCTRICA, ELECTRÓNICA Y SISTEMAS}}\\[0.3cm]
{\large \textbf{ESCUELA PROFESIONAL DE INGENIERÍA DE SISTEMAS}}\\[1cm]

\includegraphics[width=3.5cm]{logo.png}\\[1cm]

{\Large \textbf{DESARROLLO DE UN ANALIZADOR LÉXICO Y SINTÁCTICO PARA UN LENGUAJE TIPO C}}\\[1.5cm]

\begin{center}
    \textbf{PRESENTADO POR:}\\
Jhon Edy Huancuni Ccama\\[0.8cm]

\textbf{CURSO:}\\
Lenguajes de Programación\\[0.8cm]

\textbf{DOCENTE:}\\
Jose Luis Juarez Ruelas\\[0.8cm]

\textbf{SEMESTRE ACADÉMICO:}\\
2026-I
\end{center}

\vfill

{\large \textbf{PUNO -- PERÚ}}\\
{\large 2026}

\end{titlepage}

\newpage

\section*{LINKS DEL PROYECTO}
\addcontentsline{toc}{section}{LINKS DEL PROYECTO}

\begin{center}
\textbf{Link de video del proyecto}\\[0.3cm]
\url{https://drive.google.com/drive/folders/1GDHls4fVMZTqzTVsMXGGAJ9VqVojHJgL?usp=drive_link}

\vspace{1cm}

\textbf{Repositorio GitHub}\\[0.3cm]
\url{https://github.com/liamcruzticona/LP}
\end{center}

\newpage

% -------------------------
% ÍNDICE GENERAL
% -------------------------
\tableofcontents
\newpage
\listoffigures
\newpage

% -------------------------
% RESUMEN
% -------------------------
\section*{RESUMEN}
\addcontentsline{toc}{section}{RESUMEN}

El presente proyecto propone el desarrollo de un prototipo de software que realiza el análisis léxico y sintáctico de código en un lenguaje tipo C, complementado con el cálculo de entropía de Shannon para medir la complejidad informativa. Esta herramienta atiende la necesidad de los estudiantes de visualizar cómo un compilador identifica tokens, valida estructuras y cuantifica información. Se utiliza una arquitectura modular basada en Programación Orientada a Objetos y una implementación cliente-servidor con Flask en el backend y un editor web en el frontend. La metodología incluye revisión teórica, diseño de software, implementación modular, pruebas con ejemplos de código tipo C y análisis de resultados.

\textbf{Palabras clave:} análisis léxico, análisis sintáctico, compiladores, entropía de Shannon, lenguaje tipo C, ingeniería de software.

% -------------------------
% ABSTRACT
% -------------------------
\newpage
\section*{ABSTRACT}
\addcontentsline{toc}{section}{ABSTRACT}

This project proposes the development of a software prototype that performs lexical and syntactic analysis of C-type programming language code, complemented by Shannon entropy calculation to measure informational complexity. The tool addresses the need to visualize how a compiler identifies tokens, validates structures and quantifies information. A modular architecture based on Object-Oriented Programming is used, with a Flask backend and a web editor frontend. The methodology includes theoretical review, software design, modular implementation, and testing with C-type code examples.

\textbf{Keywords:} lexical analysis, syntactic analysis, compilers, Shannon entropy, C-like language, software engineering.

\newpage

% -------------------------
% PLANTEAMIENTO DEL PROBLEMA
% -------------------------
\section{PLANTEAMIENTO DEL PROBLEMA}

En la formación de estudiantes de Ingeniería de Sistemas, la comprensión de los procesos internos de un compilador es fundamental pero a menudo insuficiente. El análisis léxico y sintáctico se enseña teóricamente, pero pocas veces se presenta con una herramienta que permita ver en tiempo real cómo se tokeniza y se valida un código, generando una brecha entre teoría y práctica. Este proyecto propone cerrar esa brecha mediante un prototipo que permita ingresar código tipo C, reconocer tokens, validar estructuras sintácticas y calcular la entropía informativa del código.

% -------------------------
% JUSTIFICACIÓN
% -------------------------
\section{JUSTIFICACIÓN}

El trabajo se justifica en la necesidad de ofrecer un recurso didáctico que permita a los estudiantes comprender el análisis de compiladores de forma práctica. Al observar directamente la tokenización, la construcción del árbol sintáctico y el cálculo de la entropía, los alumnos pueden relacionar la teoría con resultados concretos. Además, el uso de programación orientada a objetos aporta una estructura modular y extensible al prototipo.

% -------------------------
% OBJETIVOS
% -------------------------
\section{OBJETIVOS}

\subsection{Objetivo general}

Evaluar la complejidad estructural e informativa de fragmentos de código tipo C mediante un prototipo orientado a objetos que realice análisis léxico, análisis sintáctico y cálculo de entropía.

\subsection{Objetivos específicos}

\begin{itemize}
    \item Diseñar un módulo de análisis léxico que identifique palabras reservadas, identificadores, literales, operadores y símbolos.
    \item Construir un módulo de análisis sintáctico que valide declaraciones, funciones y estructuras de control de un lenguaje tipo C.
    \item Implementar un módulo de cálculo de entropía que cuantifique la complejidad informativa del código fuente.
    \item Integrar los módulos en un prototipo web evaluado mediante pruebas con ejemplos de código tipo C.
\end{itemize}

% -------------------------
% METODOLOGÍA
% -------------------------
\section{METODOLOGÍA}

La metodología del proyecto se orienta a un desarrollo aplicado mediante etapas secuenciales. Se definieron componentes modulares que fueron construidos e integrados para formar un prototipo funcional.

\subsection{Diseño de arquitectura}

El sistema se compone de los siguientes elementos:

\begin{itemize}
    \item \textbf{Frontend}: editor web basado en CodeMirror, validación de entrada y visualización de resultados.
    \item \textbf{Backend}: API REST en Flask que recibe código, ejecuta análisis y devuelve resultados.
    \item \textbf{Módulos core}: clases independientes para léxico, sintaxis, tokens y entropía.
    \item \textbf{Contenedores Docker}: backend y frontend orquestados mediante Docker Compose.
\end{itemize}

\begin{figure}[H]
    \centering
    \imgstd{arqui.png}
    \caption{Arquitectura general del prototipo}
    \label{fig:arquitectura}
\end{figure}

\subsection{Pseudocódigo del flujo principal}

\begin{verbatim}
1. Recibir código fuente desde el frontend.
2. Ejecutar analizador léxico:
   - Reconocer tokens con expresiones regulares.
   - Ignorar comentarios y espacios.
3. Ejecutar analizador sintáctico:
   - Parsear declaraciones, funciones, if, while, for.
   - Manejar expresiones con precedencia de operadores.
4. Calcular entropía de Shannon:
   - Contar frecuencia de tokens.
   - Calcular probabilidades y entropía.
5. Retornar JSON con tokens, árbol sintáctico y estadísticas.
\end{verbatim}

\subsection{Entropía de Shannon}

La entropía de Shannon es una medida de la incertidumbre o complejidad informativa en un conjunto de datos, propuesta por Claude Shannon en 1948. En el contexto de este proyecto, se aplica para cuantificar la diversidad y complejidad de los tokens en el código fuente, lo que permite evaluar la riqueza léxica del programa.

La fórmula de la entropía de Shannon se define como:

\[ H = -\sum_{i=1}^{n} p_i \log_2(p_i) \]

donde:
- \( p_i \) es la probabilidad de ocurrencia del token \( i \),
- \( n \) es el número total de tipos de tokens únicos,
- \( \log_2 \) es el logaritmo en base 2.

Para calcular \( p_i \), se cuenta la frecuencia de cada token en el código y se divide por el total de tokens. Por ejemplo, si un código tiene 100 tokens y el token "int" aparece 10 veces, entonces \( p_{\text{int}} = 0.1 \).

El valor de \( H \) varía entre 0 (todos los tokens son iguales, baja complejidad) y \( \log_2(n) \) (entropía máxima, alta diversidad). En el prototipo, se normaliza la entropía dividiendo por la entropía máxima para obtener un valor entre 0 y 1, facilitando la comparación.

Esta métrica complementa el análisis léxico y sintáctico, proporcionando una dimensión cuantitativa de la complejidad del código fuente.



% -------------------------
% IMPLEMENTACIÓN
% -------------------------
\section{IMPLEMENTACIÓN}

El prototipo se implementó con los siguientes componentes:

\begin{itemize}
    \item \textbf{backend/app.py}: inicializa Flask, habilita CORS y registra el blueprint de análisis.
    \item \textbf{backend/routes/analisis.py}: define el endpoint \texttt{/analizar} y administra los errores de cada fase.
    \item \textbf{backend/core/token.py}: define la representación de tokens con tipo, valor, línea y columna.
    \item \textbf{backend/core/lexico.py}: realiza la tokenización avanzada con floats, strings, operadores compuestos y comentarios.
    \item \textbf{backend/core/sintactico.py}: parsea declaraciones, funciones, if, while, for y expresiones complejas.
    \item \textbf{backend/core/entropia.py}: calcula entropía de Shannon y estadísticas normalizadas.
    \item \textbf{frontend/index.html}: interfaz web con editor, botones de ejemplo y validación previa al análisis.
\end{itemize}

\begin{figure}[H]
    \centering
    \imgstd{frontend.png}
    \caption{Interfaz funcional}
    \label{fig:arquitectura}
\end{figure}

% -------------------------
% RESULTADOS
% -------------------------
\section{RESULTADOS}

El prototipo entrega los siguientes resultados:

\begin{itemize}
    \item Lista de tokens reconocidos por tipo, valor y posición.
    \item Árbol sintáctico en formato JSON que muestra declaraciones, expresiones y bloques.
    \item Estadísticas de entropía: valor de Shannon, entropía máxima, normalización, tipos únicos y densidad.
    \item Manejo de errores con mensajes claros y sugerencias para el usuario.
\end{itemize}

\begin{figure}[H]
    \centering
    \imgstd{json.png}
    \caption{Resultado del análisis en formato JSON}
    \label{fig:api}
\end{figure}

\subsection{Capturas de la ejecución del prototipo}

A continuación, se muestran capturas representativas del funcionamiento del sistema en diferentes etapas: interfaz, ejecución del backend, respuesta de la API y estructura del proyecto.

\begin{enumerate}
    \item Interfaz del frontend con el editor CodeMirror y un código de ejemplo cargado en el navegador en \texttt{http://localhost:8080}.

    \begin{figure}[H]
        \centering
        \imgstd{analizador.png}
        \caption{Interfaz del analizador léxico-sintáctico en el frontend}
        \label{fig:frontend}
    \end{figure}

    \item Ejecución del backend mostrando la terminal con el comando \texttt{docker-compose up --build} y el servidor Flask iniciado en el puerto 5000.

    \begin{figure}[H]
        \centering
        \imgstd{terminal.png}
        \caption{Ejecución del backend con Flask en Docker}
        \label{fig:backend_terminal}
    \end{figure}

    \item Respuesta en formato JSON de la API tras analizar un fragmento de código como \texttt{float pi = 3.14;} o estructuras condicionales como \texttt{if/else}, mostrando tokens, análisis sintáctico y entropía.

    \begin{figure}[H]
        \centering
        \imgstd{salidajson.png}
        \caption{Salida JSON del análisis léxico-sintáctico}
        \label{fig:api_json}
    \end{figure}

    \item Estructura del proyecto en VS Code, donde se observa la organización de archivos del backend y frontend.

    \begin{figure}[H]
        \centering
        \imgstd{archivos.png}
        \caption{Estructura del proyecto en el entorno de desarrollo}
        \label{fig:estructura_proyecto}
    \end{figure}

\end{enumerate}

% -------------------------
% CONCLUSIONES
% -------------------------
\section{CONCLUSIONES}

Se concluye que la construcción de un prototipo modular permite visualizar de forma práctica los procesos de análisis léxico y sintáctico de un lenguaje tipo C. El cálculo de entropía añade una dimensión cuantitativa al análisis, mostrando la complejidad informativa del código. La arquitectura orientada a objetos y la separación en módulos facilitan la comprensión del sistema y su posible ampliación.



% -------------------------
% BIBLIOGRAFÍA
% -------------------------
\section{BIBLIOGRAFÍA}

\begin{thebibliography}{9}
\bibitem{aho2012}
Aho, A. V., Lam, M. S., Sethi, R., \& Ullman, J. D. (2012). \textit{Compiladores: principios, técnicas y herramientas}. Pearson Educación.

\bibitem{cover2012}
Cover, T. M., \& Thomas, J. A. (2012). \textit{Elementos de teoría de la información}. Wiley Inter-Science.

\bibitem{pressman2015}
Pressman, R. S., \& Maxim, B. R. (2015). \textit{Ingeniería de software: un enfoque práctico}. McGraw-Hill Education.
\end{thebibliography}

\newpage

% -------------------------
% ANEXO
% -------------------------
\section{Anexo}

\subsection{Código de la arquitectura del sistema (PlantUML)}

El siguiente código corresponde al diagrama de arquitectura del sistema.

\begin{tcolorbox}[
    colback=gray!5!white,
    colframe=black,
    title=Código PlantUML - Arquitectura del sistema
]
\begin{verbatim}
@startuml

rectangle "Frontend\nCodeMirror Editor" as FE
rectangle "Backend\nFlask API" as BE

rectangle "Módulo Léxico" as LEX
rectangle "Módulo Sintáctico" as SYN
rectangle "Módulo Tokens" as TOK
rectangle "Módulo Entropía" as ENT

FE --> BE : HTTP JSON
BE --> LEX
BE --> SYN
BE --> TOK
BE --> ENT

LEX --> BE
SYN --> BE
TOK --> BE
ENT --> BE

BE --> FE : Respuesta JSON

@enduml
\end{verbatim}
\end{tcolorbox}

\end{document}
